%%%%%%%%%%%%
%
% $Autor: Wings $
% $Datum: 2019-03-05 08:03:15Z $
% $Pfad: Functions.tex $
% $Version: 4250 $
% !TeX spellcheck = en_GB/de_DE
% !TeX encoding = utf8
% !TeX root = manual 
% !TeX TXS-program:bibliography = txs:///biber
%
%%%%%%%%%%%%

\chapter{Funktionen}
In diesem Kapitel werden die Grundfunktionen des elektrischen Transportbands und deren Konfiguration erläutert.
Weiterhin werden noch zusätzliche Funktionen des Systems und ihre Verwendung dargestellt.

\section{Grundfunktionen}

\begin{itemize}
	
	\item \textbf{Betrieb mit konstanter Geschwindigkeit} \\ 
	
	Schalten Sie das System über den Hauptschalter [TeilNr01] ein und stellen Sie mithilfe des Drehreglers [TeilNr02] die gewünschte Bandgeschwindigkeit ein.
	In dieser Konfiguration versucht das Transportband die Bandgeschwindigkeit unabhängig von der Belastung beizubehalten.\\
	Es können Bandgeschwindigkeiten zwischen 1,8 und 7,6 cm/s eingestellt werden.
	
	\item \textbf{Umkehr der Laufrichtung}\\
	Das Transportband kann in zwei Laufrichtungen betrieben werden.\\
	Stellen Sie vor Umkehr der Laufrichtung sicher, dass das Transportband vollständig stillsteht.\\
	Dazu sollte sich der Drehregler [TeilNr02] auf der niedrigsten Einstellung [0] befinden oder der Hauptschalter [TeilNr01] auf der Position [O] stehen.
	\\ 
	Stellen Sie nun den Richtungswahlschalter [TeilNr03] von Position01 auf Position02.
	Das System kann anschließend über den Hauptschalter [TeilNr01] eingeschaltet werden.
	Die gewünschte Bandgeschwindigkeit kann erneut mithilfe des Drehreglers [TeilNr02]  eingestellt werden.
	
\end{itemize}

\section{Erweiterte Funktionen}

\begin{itemize}
	
	\item \textbf{Betrieb mit Lichtschranken}\\
	TEXT
	
	\item \textbf{Betrieb mit programmiertem Bewegungsablauf} \\ 
	TEXT
	
\end{itemize}
